% =============================================================================
% SECTION 10: JUSTIFICATION OF CONSTRAINTS IMPOSED BY THE DESIGN
% =============================================================================
\section{Justification of Constraints Imposed by the Design}
\label{sec:constraints}

This section identifies and justifies the constraints that the Space Segment design imposes on production, operation, transportation/storage, and maintainability.

% -----------------------------------------------------------------------------
\subsection{Production}
\label{sec:constraints_production}

\subsubsection{Serial Production Scale and Rate}

The constellation design (72~\texttt{SATL} + 6~\texttt{SATG}) imposes a serial production constraint unprecedented for European spacecraft programmes:
\begin{itemize}
    \item \textbf{SATL production:} 72 flight units organised in 8 plane-lots of 9 spacecraft each. The production facility must sustain a throughput of $\ge 2$ completed \texttt{SATL} units per month to meet the launch manifest timeline, requiring parallel AIT stations and batch-tested component inventories.
    \item \textbf{SATG production:} 6 flight units manufactured individually with full acceptance testing per unit.
    \item \textbf{Component procurement:} The ``Careful COTS'' philosophy for \texttt{SATL} (\textbf{AD-03}, Section 7.2) mandates lot-acceptance radiation testing of every automotive/industrial-grade processor and FPGA lot under the project radiation-assurance rules \emph{prior to} board assembly. Lot rejection rates and lead-time variability must be absorbed by the procurement schedule with $\ge 10\%$ component over-procurement margin.
\end{itemize}

\subsubsection{Cleanroom and Contamination Control}

\begin{itemize}
    \item \textbf{Payload and OISL Optical Assemblies (\texttt{.04}, \texttt{.06.02}):} All cooled SWIR/MWIR telescope assemblies, dichroic beam-splitters, cryogenic HgCdTe FPA modules, and 1,550~nm laser OISL optical heads must be integrated and maintained under \textbf{ISO Class~5 (Class~100)} cleanroom conditions with continuous high-purity dry $\mathrm{GN_2}$ purging. Optical surfaces require sealed aperture covers whenever not in active test or calibration. Material selection is constrained to $\mathrm{CVCM} < 0.01\%$ and $\mathrm{TML} < 1.0\%$ to prevent molecular condensation on cryogenic optics ($65$--$95$~K FPA, $180$--$210$~K mirrors) that would absorb in the $2.7~\mu\mathrm{m}$ (SWIR) and $4.2$--$4.4~\mu\mathrm{m}$ (MWIR) science bands (\texttt{R-009}, \texttt{R-025}).
    \item \textbf{Platform Bus Integration:} Conducted in \textbf{ISO Class~8 (Class~100,000)} facilities.
    \item \textbf{Propulsion System:} Propellant loading and pressurant charging are performed in Hazardous Processing Facilities at the launch site, not at the AIT facility, to avoid contamination-hazard co-location.
\end{itemize}

\subsubsection{Calibration Constraints}

\begin{itemize}
    \item Every \texttt{SATL} and \texttt{SATG} flight unit requires individual radiometric calibration of both SWIR ($2.65$--$2.85~\mu\mathrm{m}$) and MWIR ($4.15$--$4.45~\mu\mathrm{m}$) FPA channels at their cryogenic operating temperature ($65$--$95$~K) inside a thermal-vacuum chamber with an optical port and calibrated blackbody sources. This constrains TVAC chamber scheduling.
    \item Geometric distortion mapping of the NFOV ($1.5^\circ \times 1.5^\circ$) must achieve pixel-level co-registration between SWIR and MWIR channels to validate two-colour ratioing performance (\texttt{R-019}, \texttt{R-022}).
    \item Each OISL terminal (\texttt{.06.02}) requires optical acquisition and bit-error-rate (BER) ground testing at representative free-space-equivalent distances using collimated beam simulators.
\end{itemize}

\subsubsection{Recurring Production Acceptance Constraints}

Per the Verification Plan (\textbf{AD-04}), every recurring \texttt{SATL} FM undergoes:
\begin{itemize}
    \item Workmanship random vibration (acceptance spectrum, reduced from qualification levels).
    \item Shortened TVAC acceptance screening (reduced cycle count relative to QM/PFM).
    \item Radiometric and geometric payload calibration.
    \item Functional chain test with flight software, including FDIR fault injection.
    \item Mass properties measurement (weigh and CG) for comparison against the mass budget.
\end{itemize}
The 72-unit SATL fleet implies $> 500$ individual environmental test runs (vibration + TVAC + EMC), constraining test facility availability across the production timeline.

% -----------------------------------------------------------------------------
\subsection{Operation}
\label{sec:constraints_operation}

\subsubsection{Pointing, Acquisition, and Tracking (PAT) Constraints}

\begin{itemize}
    \item \textbf{LEO--GEO Cross-Layer OISL PAT:} The 1,550~nm laser OISL terminals must acquire and maintain optical lock over the full $34{,}800$--$45{,}000$~km range (\texttt{R-033}) despite relative orbital motion between the 7.350~km/s LEO node and the near-stationary GEO node. Beam pointing must be maintained below $< 1.5~\mu\mathrm{rad}$ while the \texttt{SATL} bus traverses ${\sim}13.7$~orbits/day. During Mode~N2 (Agile Stereo Tracking), the OISL terminals must sustain lock \emph{concurrently} with independent payload gimbal slewing ($\ge 15^\circ$/s), which induces structural micro-vibration coupling into the OISL fine-steering mirror control loop.
    \item \textbf{LEO--LEO Intra-/Cross-Plane OISL PAT:} Constant 5,047~km chord with zero relative azimuth rate for intra-plane links ensures straightforward PAT lock. Cross-plane links (5,647--6,144~km at equator, dynamically varying geometry toward $\pm 80^\circ$ latitude) require agile coarse-pointing periscopes with continuous handover logic at the polar cross-seam to sustain the low-latency inter-satellite mesh.
\end{itemize}

\subsubsection{Thermal and Power Constraints During Continuous Dual-Band Streaming}

\begin{itemize}
    \item \textbf{Cryocooler Duty Cycle:} The pulse-tube cryocoolers maintaining HgCdTe FPAs at $65$--$95$~K ($\pm 0.1$~K) operate \emph{continuously} during all modes (N1 through N5). Each cryocooler dissipates TBD~W of waste heat into dedicated cryocooler radiator panels, which must maintain cold-sink temperature below the cryocooler reject temperature at all orbital beta angles, including worst-case $\beta_\odot \approx 0^\circ$ with maximum Earth albedo ($0.38$) and IR ($260$~W/m$^2$) flux at 1,000~km.
    \item \textbf{Mode~N2 Peak Power:} During the $\le 25$--$30$~s Mode~N2 active tracking window, all four OISL terminals, the payload processor, the gimbal/FSM drive electronics, and the SWIR/MWIR FPAs operate simultaneously at full rate ($30$~Hz, $640 \times 640$ pixels, dual-band). This imposes a transient peak power demand on the EPS that must remain within the $\ge 550$~W (SATL) / $\ge 4{,}200$~W (SATG) EOL average power budget with $\ge 15\%$ positive margin. Battery DoD during eclipsed Mode~N2 events must remain $\le 30\%$ (SATL) / $\le 65\%$ (SATG).
    \item \textbf{Data Handling Throughput:} The OBDH payload co-processor (FPGA) must sustain real-time proton despiking, two-colour ratioing, and spatio-temporal clutter rejection on the full dual-band $640 \times 640 \times 30$~Hz data stream ($\ge 196.6$~Mbps raw at 8-bit) without frame drops. Simultaneously, it must format and encrypt CCSDS packets and route them to the OISL or S-band TT\&C transmitters within the Step~1 ($\le 10$~s, \texttt{R-042}) and Step~3 ($\le 25$--$30$~s) latency allocations.
\end{itemize}

\subsubsection{CONOPS and Coverage Constraints}

\begin{itemize}
    \item \textbf{Staggered Maintenance Burns:} Orbit maintenance and CAM are constrained to at most one \texttt{SATL} node per coverage region in manoeuvre mode at any time (\texttt{R-008}, \texttt{R-015}). The Ground Segment flight dynamics system must deconflict station-keeping windows across all 72 SATL nodes and 6 SATG nodes to preserve $\ge 5$ simultaneous LEO nodes visible over the Baltic ($53^\circ$--$66^\circ$N) at all times.
    \item \textbf{TT\&C Contact Windows:} S-band TT\&C for \texttt{SATL} nodes is constrained by the 1,000~km orbit period ($T = 105.12$~min) and the geographic distribution of the Ground Station Network (\abbr{WBS} \texttt{123456.02.05}). Each \texttt{SATL} node has ${\sim}8$--$12$~min of direct S-band ground contact per pass. Outside direct contact windows, all TC must be routed via the OISL mesh relay network (Mode~N3, \texttt{R-044}). \texttt{SATG} nodes at 35,786~km have near-continuous S-band and Ka-band contact with European ground stations.
    \item \textbf{SAA Operational Degradation:} During SAA transits (${\sim}8$--$12$~min, 4--6 times/day per SATL node), the infrared FPAs experience elevated proton-induced transient pixel spikes requiring real-time despiking (\texttt{R-019}). Detection sensitivity ($P_d$) may be temporarily reduced during SAA passage; the warning chain availability analysis accounts for this by ensuring $\ge 5$--$7$ non-SAA SATL nodes remain visible over the Baltic at all times.
\end{itemize}

\subsubsection{OISL Mesh Relay Constraints}

\begin{itemize}
    \item The mandatory OISL relay (Mode~N3) imposes a continuous $\ge 500$~Mbps per-link capacity reservation on all 4 SATL terminals and cross-layer GEO--LEO trunks, consuming electrical power and thermal dissipation budgets at all times (not only during events).
    \item Polar cross-seam handover ($\pm 80^\circ$ latitude) requires dynamic re-acquisition of cross-plane OISL links within $< 3$~s to remain within the Step~2 latency allocation (\texttt{R-044}).
    \item COMSEC hardware encryption modules (\texttt{R-031}) impose a fixed processing latency and power overhead on every OISL data packet, independent of link traffic volume.
\end{itemize}

% -----------------------------------------------------------------------------
\subsection{Transportation and Storage}
\label{sec:constraints_transport}

\subsubsection{Containerized Transportation}

\begin{itemize}
    \item Spacecraft are transported in hermetically sealed, climate-controlled shipping containers: $20^\circ\mathrm{C} \pm 3^\circ\mathrm{C}$, relative humidity $< 40\%$, continuous $\mathrm{GN_2}$ purge of optical apertures.
    \item Damped shock-isolation frames limit quasi-static and dynamic transport loads to $\le 3.0\,g$ along all axes, monitored by continuous triaxial shock/temperature data loggers.
    \item Container dimensions must accommodate the stowed \texttt{SATL} envelope (including folded solar arrays and stowed antennas) and must interface with standard air-cargo (C-130, An-124) and road-transport platforms.
    \item \texttt{SATG} spacecraft ($2{,}400$--$3{,}100$~kg wet) require heavy-lift transport and dedicated handling frames compatible with launch-site clean-room crane capacities.
\end{itemize}

\subsubsection{Launch Vehicle Interface Constraints}

\begin{itemize}
    \item \textbf{SATL:} The 9-satellite cylindrical multi-satellite dispenser ring stack (${\sim}2{,}300$--$2{,}800$~kg total) must fit within the launcher fairing usable dynamic envelope. Individual SATL structural interfaces are constrained by the dispenser adapter ring bolt pattern, separation spring force envelope ($\Delta v_{\text{sep}} \approx 0.5$--$1.5$~m/s), and stowed stiffness requirements ($f_{\text{axial}} \ge 35$~Hz, $f_{\text{lateral}} \ge 20$~Hz).
    \item \textbf{SATG:} Single- or dual-manifest heavy-lift vehicle interface. The launch adapter and clamp-band interface must accommodate the $\le 3{,}100$~kg wet mass with stowed stiffness $f_{\text{axial}} \ge 31$~Hz, $f_{\text{lateral}} \ge 15$~Hz. Separation shock $1{,}500$--$2{,}000\,g$ SRS above 1,000~Hz.
    \item \textbf{Launch Safety Inhibits:} During fairing encapsulation through ascent, all propulsion valves, high-power RF transmitters, 1,550~nm laser OISL diodes, and pulse-tube cryocoolers must remain under $\ge 3$ independent electrical or mechanical safety inhibits. HDRM and aperture covers remain mechanically locked.
\end{itemize}

\subsubsection{Pre-Launch Storage and Shelf-Life Constraints}

\begin{itemize}
    \item \textbf{Maximum Storage Duration:} Completed flight units may be stored for up to 24~months in controlled cleanroom conditions (ISO~8 for bus, ISO~5 for optics) with periodic powered health checks.
    \item \textbf{Li-ion Battery Storage:} Battery state-of-charge maintained at $40$--$50\%$ storage charge to minimise calendar ageing and capacity fade. Periodic top-off/equalization cycles every 3--6~months.
    \item \textbf{HgCdTe FPA Detector Shelf Life:} Cryogenic FPA modules stored at ambient temperature in sealed $\mathrm{GN_2}$-purged containers. No shelf-life degradation mechanism at room temperature; dark-current baseline must be re-verified if storage exceeds 18~months.
    \item \textbf{Propellant Tanks:} Propulsion subsystem delivered dry to launch site. Propellant loading and pressurant charging performed at the Hazardous Processing Facility during the launch campaign. Tank proof-pressure and leak-test certificates are valid for TBD~months from acceptance.
    \item \textbf{OISL Laser Diode Ageing:} 1,550~nm laser diodes stored unpowered at $20^\circ\mathrm{C} \pm 3^\circ\mathrm{C}$; no significant shelf-life degradation at room temperature. Pre-launch re-verification of output power and wavelength stability if storage exceeds 24~months.
\end{itemize}

% -----------------------------------------------------------------------------
\subsection{Maintainability}
\label{sec:constraints_maintainability}

\subsubsection{On-Orbit Software and Firmware Maintainability}

\begin{itemize}
    \item \textbf{Flight Software Patching (Mode~N5, \texttt{R-004}, \texttt{R-041}):} The OBDH flight software and payload co-processor FPGA bitstreams are designed for authenticated over-the-air uplink of refined plume-detection algorithms, updated two-colour ratioing filter coefficients, and security patches derived from the 5-year Ground Segment operational archive. All uploads are protected by COMSEC hardware encryption (\texttt{R-031}) and require dual-authenticated Ground Segment operator authorization.
    \item \textbf{FPGA Reconfiguration:} SATL payload co-processor FPGAs support partial reconfiguration to allow in-orbit updates to proton despiking, clutter rejection, and spectral ratioing kernels without full FPGA re-flash. SATG QML-V rad-hard FPGAs require full re-flash via protected configuration memory banks.
    \item \textbf{COMSEC Key Rotation (\texttt{R-031}):} Cryptographic key material is rotated autonomously on schedule and can be emergency-rekeyed via authenticated uplink. Key storage is in tamper-resistant hardware security modules.
\end{itemize}

\subsubsection{FDIR and Redundancy Switching}

\begin{itemize}
    \item \textbf{Dual Cold/Warm Block Redundancy (\texttt{R-017}):} All critical bus chains (OBDH, ADCS, EPS, OISL transceivers, cryocooler drive electronics) implement A/B redundant unit pairs with autonomous Level~2 FDIR failover ($< 2.0$~s). After failover, the faulted A-side unit is electrically isolated; ground operators can command a reversion attempt during the next scheduled maintenance window.
    \item \textbf{GEO Relay Chain Single-Point Failures:} The GEO--LEO optical relay chain is a critical path in the warning latency budget (Step~2, $\le 3.0$~s OISL routing, \texttt{R-044}). Loss of a single \texttt{SATG} OISL relay node is mitigated by the $60^\circ$-spaced 6-node GEO ring: routing tables dynamically re-converge within $< 3$~s to bypass the faulted GEO node and use an adjacent \texttt{SATG} as the relay gateway. Within each \texttt{SATG}, the cross-layer OISL terminal is backed by a redundant B-side transceiver (\texttt{R-017}).
    \item \textbf{Cryocooler Redundancy:} Loss of a single pulse-tube cryocooler renders both SWIR and MWIR FPAs non-operational on the affected spacecraft (they share a common cold-finger thermal bus). Cryocooler drive electronics are dual-redundant; a second cryocooler compressor head is TBD at current design maturity. The constellation design tolerates the loss of individual SATL nodes (5--7 visible over Baltic at all times).
\end{itemize}

\subsubsection{Pre-Launch Ground Accessibility and LRU Replaceability}

\begin{itemize}
    \item \textbf{Line-Replaceable Unit (LRU) Access:} The SATL and SATG structural design provides panel-removable access to the following LRUs during AIT without disturbing adjacent subsystems:
    \begin{itemize}
        \item OBDH processor boards (dual A/B units)
        \item ADCS reaction wheel assemblies and star tracker heads
        \item EPS battery modules and PCDU
        \item OISL transceiver modules (\texttt{.06.02})
        \item S-band TT\&C transponder (\texttt{.06.01})
        \item Cryocooler compressor and drive electronics
    \end{itemize}
    Any LRU replacement performed after spacecraft-level environmental testing requires a re-test at acceptance levels (workmanship vibration + abbreviated TVAC) on the affected subsystem interface per the Verification Plan (\textbf{AD-04}).
    \item \textbf{Payload Non-Replaceability:} The integrated SWIR/MWIR telescope assembly with dichroic beam-splitter and cryogenic FPA cold finger is \emph{not} an LRU. It is permanently integrated during payload-level AIT and cannot be replaced at spacecraft level without full disassembly and optical re-alignment. Any payload-level failure after spacecraft integration requires return to the payload AIT facility.
    \item \textbf{Propulsion Non-Replaceability:} Once propulsion tanks and feed lines are proof-tested and welded at subsystem level, they are permanently integrated. Post-integration propulsion anomalies require spacecraft-level disposition.
\end{itemize}
